\chapter{Data Acquisition, Feature Engineering, and Modelling}

This chapter describes how gesture data is collected, how it is represented,
and how the recognition model is trained and deployed. The feature-%
representation study in Section~\ref{sec:features} is the central technical
contribution of the project.

\section{Data Acquisition}

Data is collected with the master in CSV mode (Chapter~4) connected to a
laptop, driven by a Python recorder. The operator selects a label, the recorder
shows a short countdown, and a fixed-length window is captured while the
gesture is performed. Each saved sample is a CSV file of 75 rows
($3\,\text{s} \times 25\,\text{Hz}$) and 45 channels. The recorder forward-fills
the most recent packet from each finger onto a uniform time grid and flags
samples in which any finger's data went stale, so that dropped-packet windows
can be discarded.

\section{Signal Window and Channels}

A gesture is represented by a $75 \times 45$ window: 75 time steps of the 45
channels defined in Chapter~4. This window length and channel order are fixed
across data collection, training, and on-device inference. Matching them
exactly—\emph{train/inference parity}—is essential; a mismatch in window
length, sampling rate, or channel order silently degrades or breaks live
accuracy even when offline metrics look perfect.

\section{Feature Representation}\label{sec:features}

The most important design question proved to be \emph{what} the model sees, not
which classifier is used.

\subsection{The Problem with Absolute Orientation for Static Handshapes}

A natural first representation is the absolute fused orientation of each finger.
This failed in live use for static fingerspelling in two distinct ways. First,
the absolute pitch/roll of each finger is measured relative to the start-up
calibration pose, which a user cannot reproduce exactly between sessions; a few
degrees of difference in how the hand is held shifts every channel and causes
confident misclassification. Second, even with perfect calibration, letters
that differ only by small thumb placement (e.g.\ \textbf{A} vs \textbf{T}) lie
within the angular noise floor of the sensors and cannot be separated.

\subsection{Relative (Orientation-Invariant) Features}

To address the calibration sensitivity, an orientation-invariant transform was
introduced. For the orientation-bearing channels ($P,R,Y,a_x,a_y,a_z$), the
per-time-step mean across the five fingers is subtracted, leaving only the
\emph{inter-finger} differences; the angular-rate channels are left unchanged.
Because a global rotation of the hand (or a calibration offset) adds the same
amount to all five fingers, subtracting their mean cancels it, while the
relative finger configuration—which defines the handshape—is preserved. This
made static-letter features invariant to global hand tilt and to the
calibration pose.

Relative features removed the calibration sensitivity but, as expected, could
not overcome the sensor-resolution limit: the thumb-cluster letters remained
confusable, because the information that separates them was never captured by
the sensors in the first place.

\subsection{Dynamic Gestures with Absolute Features}

The resolution of both problems was to change the gestures rather than only the
features. Each letter (except the inherently dynamic \textbf{J} and \textbf{Z})
was re-recorded as a short, distinctive \emph{motion} that engages specific
fingers—for example a thumb fold, an index-and-thumb tap, or a thumb-and-pinky
meeting—so that the discriminative information lives in the angular-rate
($g$) channels. The angular rate is a measurement of motion and is inherently
independent of the start-up calibration, so \emph{absolute} features (no
mean-centering) can be used. This dynamic-gesture, absolute-feature formulation
is used for both the letter and the word models. A small number of practical
rules emerged: the motion must fill the whole capture window (a gesture with a
long still tail collapses toward a low-motion ``attractor'' class), and within
a confusable group the motions must differ clearly in axis, direction, or
repetition count.

The same logic was confirmed for word signs. Word pairs distinguished only by
global hand orientation or translation (e.g.\ \textsc{you} vs \textsc{me})
collapsed under mean-centering—which deliberately discards exactly that
information—and were correctly separated once absolute features were used.

\section{Augmentation and Dataset Splitting}

Recording every gesture across many sessions is time-consuming, so synthetic
augmentation is used to improve robustness to session-to-session variation.
For each training sample, several copies are generated with small random
per-finger angle offsets, small per-finger gyroscope-bias offsets, and
per-reading sensor noise—simulating realistic calibration and thermal variance
without re-recording. Critically, the train/test split is performed
\emph{before} augmentation and on the original samples only, so that no
augmented copy of a test sample can appear in training. This leakage-safe
procedure keeps the held-out test set an honest measure of generalisation.

\section{Network Architecture}

The classifier is a compact one-dimensional convolutional neural network
\cite{ronao2016}. The flattened input is reshaped to $75 \times 45$ and passed
through two convolutional blocks—each a 1-D convolution, batch normalisation,
max-pooling, and dropout (16 then 32, or 32 then 64, filters, kernel size
3)—followed by a flatten, a dense layer, dropout, and a softmax output over the
classes. The architecture was authored in Keras through Edge Impulse's expert
mode \cite{edgeimpulse}. Training uses categorical cross-entropy with the Adam
optimiser, a 20\% validation split, and standard-scaler normalisation of the
inputs, whose per-channel means and standard deviations are exported so the
phone can reproduce the identical normalisation at inference time.

\section{Quantisation and Deployment}

The trained model is exported as an 8-bit-integer (int8) quantised TensorFlow
Lite file \cite{tflite}. Quantisation reduces the model size roughly four-fold
and speeds up inference, and on this data it incurred no measurable accuracy
loss. The deployed model occupies about 137\,KB of flash and 22.7\,KB of RAM
and runs in approximately 1\,ms per inference (Chapter~7). Because the model's
input and output tensors are both int8, the application must quantise the
normalised input and dequantise the output using the scale and zero-point
published in the model—handled in the mobile application (Chapter~6).
